% !PRINTED 686
\weberpage{686}
%<<CONT
eine Einheitswurzel ist. Das Exponentensystem der aus $\alpha$ abgeleiteten
 reducirten Zahl $\alpha_0$ ist durch $\alpha$ selbst völlig bestimmt.
Hierdurch ist der Satz bewiesen:

II. \gesperrt{Zu jeder (von Null verschiedenen) Zahl $\alpha$ des
Körpers $\Omega$ giebt es $w$ und nicht mehr reducirte
Zahlen $\varrho\,\alpha_0$.}

Wenn statt des Fundamentalsystemes $(\varepsilon_1, \varepsilon_2, \ldots, \varepsilon_{\nu-1})$ ein
anderes angewendet wird, so erleiden die Exponenten $\xi_1, \xi_2, \ldots, \xi_{\nu-1}$
eine ganzzahlige lineare Substitution von der Determinante $\pm 1$.
Eine reducirte Zahl kann dann aufhören, für das neue Fundamentalsystem
 reducirt zu sein. Der Begriff der reducirten Zahl
ist also von der Wahl des Fundamentalsystemes abhängig. Die
Anzahl der aus einer gegebenen Zahl abgeleiteten reducirten
Zahlen ist aber immer dieselbe.

Durch diesen Satz haben wir den Zweck erreicht, aus
dem ganzen unendlichen Systeme der unter einander associirten
Zahlen eine bestimmte endliche Anzahl von Repräsentanten
herausgehoben zu haben.

\begin{center}
\S.\,189.\\[4pt]
\gesperrt{Grenzen der Anzahl der durch ein Ideal theilbaren
ganzen Zahlen des Körpers $\Omega$.}
\end{center}
% [source note] The following finiteness assertion and definition of T omit counting modulo association. Infinitely many units already have absolute norm 1. The intended finite count is of non-associated integers (equivalently the corresponding principal ideals), explicitly specified on printed p. 690. The source also changes the endpoint from norm at most t here to natural numbers less than t below; this can alter a boundary term but not the limiting quotient. The wording is preserved.

Unsere früheren Betrachtungen haben ergeben, dass jede
ganze Zahl eines Körpers $\Omega$ nur eine endliche Anzahl von Idealen
zu Theilern hat. Da jedes ganze Ideal ein Theiler seiner Norm
ist, so giebt es also auch nur eine endliche Anzahl von ganzen
Zahlen in $\Omega$, deren absolute Norm eine gegebene Grenze nicht
übersteigt.

Aus diesen Zahlen greifen wir jetzt wieder einen Theil
heraus, und fragen nach der Anzahl $T$ aller \gesperrt{ganzen Zahlen}
des Körpers $\Omega$, deren absolute Norm eine positive Grösse $t$
nicht überschreitet, und die durch ein gegebenes Ideal $\mathfrak{a}$ theilbar
 sind.

Nehmen wir als vorläufiges Beispiel den rationalen Körper,
so ist dort $T$ die Anzahl der natürlichen Zahlen, die kleiner als
$t$ und durch eine gegebene ganze Zahl $m$ theilbar sind, also, in
der früher (\S.\,182) gebrauchten Bezeichnung $T = E\left(\frac{t}{m}\right)$.
